\chapter{Réalisation}

\section{Introduction}
Ce chapitre détaille l'implémentation technique de la plateforme eduGenius. Nous y décrivons l'environnement de développement, l'intégration des API d'IA, et la stratégie de validation du logiciel.

\section{Environnement de Développement}
\subsection{Outils et Logiciels}
\begin{itemize}
    \item \textbf{Backend} : Node.js avec le framework Express.
    \item \textbf{Frontend} : Next.js 16 avec TypeScript.
    \item \textbf{Persistence} : MongoDB Atlas (Cloud).
    \item \textbf{Stockage} : Bucket AWS S3.
    \item \textbf{IA} : Google Generative AI SDK (Gemini).
\end{itemize}

\section{Implémentation des Services Clés}

\subsection{Gestion de l'État Global (Frontend)}
Pour centraliser les informations de l'utilisateur connecté et son rôle, nous avons utilisé l'API \textbf{Context} de React. Le \texttt{AuthContext} permet de diffuser l'état d'authentification à travers toute l'application sans "Prop Drilling".

\begin{lstlisting}[language=JavaScript, caption=Exemple de hook personnalisé useAuth]
export const AuthProvider = ({ children }) => {
    const [user, setUser] = useState(null);
    const [loading, setLoading] = useState(true);

    useEffect(() => {
        const token = localStorage.getItem('token');
        if (token) {
            checkAuth(token).then(u => setUser(u)).finally(() => setLoading(false));
        }
    }, []);

    return <AuthContext.Provider value={{ user, setUser }}>{children}</AuthContext.Provider>;
};
\end{lstlisting}

\subsection{Standardisation de l'Interface Utilisateur (Architecture Orientée Composants)}
Pour garantir une expérience utilisateur fluide et maintenir une cohérence visuelle à travers toute la plateforme, nous avons adopté une architecture orientée composants (Component-Based UI). L'un des piliers de cette approche a été la création de composants structurels réutilisables, tels que le composant \texttt{PageHeader}.

Ce composant centralise les éléments de navigation, les animations (via \texttt{framer-motion}) et la typographie, remplaçant ainsi le code dupliqué sur plus de 10 vues du tableau de bord étudiant (cours, évaluations, messagerie, performances, etc.). Cette refactorisation a permis non seulement d'améliorer la maintenabilité du code, mais aussi d'assurer une esthétique uniforme sur l'ensemble des modules de l'application.

\subsection{Middleware de Sécurité (Backend)}
Chaque requête vers l'API est filtrée par un middleware de protection qui décode le token JWT et injecte l'utilisateur dans l'objet \texttt{req}.

\begin{lstlisting}[language=JavaScript, caption=Middleware de protection des routes]
const protect = async (req, res, next) => {
    let token = req.headers.authorization?.split(' ')[1];
    if (!token) return res.status(401).json({ message: 'Non autorise' });

    try {
        const decoded = jwt.verify(token, process.env.JWT_SECRET);
        req.user = await User.findById(decoded.id).select('-password');
        next();
    } catch (err) {
        res.status(401).json({ message: 'Token invalide' });
    }
};
\end{lstlisting}

\section{Implémentation du Cœur d'IA (Gemma 3:4b via Ollama)}
L'innovation majeure réside dans le service \texttt{OllamaService.js}. Nous avons privilégié l'utilisation d'une instance Dockerisée d'Ollama.

\begin{lstlisting}[language=JavaScript, caption=Orchestration de l'IA locale Gemma 3]
async generateSmartContent(text, mode = 'summary') {
    const prompt = mode === 'summary' 
        ? `Resume ce cours de maniere FUN : ${text}`
        : `Genere 5 questions QCM en JSON : ${text}`;

    const response = await axios.post('http://localhost:11434/api/generate', {
        model: 'gemma3:4b',
        prompt: prompt,
        format: 'json',
        stream: false
    });
    return response.data.response;
}
\end{lstlisting}

\subsection{Configuration Matérielle Requise}
Pour une fluidité optimale lors des pics d'utilisation, le serveur de déploiement chez BES2I a été configuré avec :
\begin{itemize}
    \item \textbf{CPU} : Intel Xeon avec 16 Threads (pour les calculs sans GPU).
    \item \textbf{GPU} : NVIDIA RTX 4060 Ti (8GB VRAM) permettant une inférence en temps réel.
    \item \textbf{RAM} : 32 GB DDR5.
\end{itemize}

\subsection{Gestion de la Messagerie en Temps Réel}
Pour simuler l'interactivité de Microsoft Teams, nous avons utilisé Socket.io pour diffuser les messages instantanément sans recharger la page.

\begin{lstlisting}[language=JavaScript, caption=Gestionnaire de messages Socket.io]
socket.on('send-message', async ({ roomId, content, senderId }) => {
    try {
        const msg = await Message.create({
            senderId,
            content,
            classId: roomId,
            messageType: 'class'
        });
        // Diffusion uniquement aux membres de la room 'chat:classId'
        io.to(`chat:${roomId}`).emit('new-message', msg);
    } catch (err) {
        socket.emit('message-error', { error: 'Echec de l envoi' });
    }
});
\end{lstlisting}

\section{Sécurité et Authentification}
Nous avons implémenté un système de "Double Middleware" pour sécuriser les routes sensibles.
\begin{itemize}
    \item \textbf{authenticate} : Vérifie que le token JWT fourni dans l'en-tête \texttt{Authorization} est valide.
    \item \textbf{authorize} : Vérifie que le rôle de l'utilisateur (ex: 'admin') a les permissions nécessaires pour l'action demandée.
\end{itemize}

\section{Environnement de Déploiement}
Pour assurer une disponibilité continue et une scalabilité aisée, nous avons adopté une stratégie de déploiement multi-cloud.

\begin{table}[H]
\centering
\begin{tabular}{|l|l|p{8cm}|}
\hline
\textbf{Composant} & \textbf{Hébergeur} & \textbf{Raison du Choix} \\ \hline
Frontend (Next.js) & Vercel & Optimisation native pour Next.js, Edge Functions, et CI/CD intégré. \\ \hline
Backend (Node.js) & Render & Support natif de Docker, gestion facile des variables d'environnement. \\ \hline
Base de Données & MongoDB Atlas & Base de données managée, sauvegardes automatiques et scaling horizontal. \\ \hline
Fichiers (PDF) & AWS S3 & Stockage objet hautement disponible et sécurisé. \\ \hline
\end{tabular}
\caption{Environnements de déploiement d'eduGenius}
\end{table}

\section{Implémentation de la Logique Frontend}
\subsection{Hooks Personnalisés}
Nous avons développé plusieurs hooks personnalisés pour encapsuler la logique métier complexe et favoriser la réutilisation du code.

\begin{lstlisting}[language=JavaScript, caption=Hook personnalisé useSocket pour la gestion du temps réel]
export const useSocket = (roomId) => {
    const [messages, setMessages] = useState([]);
    const socketRef = useRef();

    useEffect(() => {
        socketRef.current = io(process.env.NEXT_PUBLIC_SOCKET_URL);
        socketRef.current.emit('join-room', roomId);

        socketRef.current.on('new-message', (msg) => {
            setMessages((prev) => [...prev, msg]);
        });

        return () => socketRef.current.disconnect();
    }, [roomId]);

    const sendMessage = (content) => {
        socketRef.current.emit('send-message', { roomId, content });
    };

    return { messages, sendMessage };
};
\end{lstlisting}
\section{Qualité du Code et Bonnes Pratiques}
Pour garantir la maintenabilité et la robustesse d'eduGenius, nous avons appliqué plusieurs standards de l'industrie.

\subsection{Clean Architecture et Séparation des Préoccupations}
Nous avons suivi les principes de la \textbf{Clean Architecture} en séparant strictement :
\begin{itemize}
    \item \textbf{Les Entités} : Modèles Mongoose définissant la structure pure des données.
    \item \textbf{Les Repositories} : Couche d'accès aux données isolant les requêtes MongoDB.
    \item \textbf{Les Services} : Logique métier pure (ex: calcul des moyennes, orchestration IA).
    \item \textbf{Les Contrôleurs} : Gestion des entrées/sorties HTTP.
\end{itemize}

\subsection{Linters et Formatage}
L'utilisation d'\textbf{ESLint} et de \textbf{Prettier} a permis de maintenir une base de code homogène, facilitant ainsi les revues de code et la collaboration future.

\section{Tests et Validation}
\subsection{Tests d'Acceptation Utilisateur (UAT)}
Un panel de 5 enseignants et 10 étudiants de l'EPI a testé la plateforme en conditions réelles.

\begin{table}[H]
\centering
\begin{tabular}{|l|c|p{8cm}|}
\hline
\textbf{Critère} & \textbf{Note (/5)} & \textbf{Commentaire} \\ \hline
Facilité d'usage & 4.5 & Interface intuitive, proche de Teams. \\ \hline
Vitesse de l'IA & 3.8 & Parfois lent pour les PDF très longs (>50 pages). \\ \hline
Qualité des résumés & 4.2 & Ton divertissant très apprécié par les étudiants. \\ \hline
Stabilité Vidéo & 4.7 & Excellente fluidité même en 4G. \\ \hline
\end{tabular}
\caption{Résultats des tests d'acceptation utilisateur (UAT)}
\end{table}

\subsection{Plan de Tests Techniques}
Le tableau suivant récapitule notre plan de test pour garantir la robustesse d'eduGenius.

\begin{table}[H]
\centering
\begin{tabular}{|l|p{6cm}|p{6cm}|}
\hline
\textbf{Type de Test} & \textbf{Objectif} & \textbf{Résultat} \\ \hline
Tests Unitaires & Validation des fonctions de calcul de moyenne. & 100\% Succès. \\ \hline
Tests d'Intégration & Communication entre le Frontend et l'API S3. & Upload/Download fonctionnel. \\ \hline
Tests de Charge & 50 utilisateurs simultanés sur le chat. & Latence moyenne < 80ms. \\ \hline
Tests d'Accessibilité & Conformité au contraste et lecture d'écran. & Score Lighthouse : 94/100. \\ \hline
\end{tabular}
\caption{Plan de tests et résultats de validation}
\end{table}

\section{Déploiement et CI/CD}
La plateforme a été déployée sur une infrastructure Cloud hybride :
\begin{itemize}
    \item \textbf{Frontend} : Déployé sur \textbf{Vercel} pour bénéficier de l'optimisation des images et du CDN global.
    \item \textbf{Backend} : Hébergé sur \textbf{Render} avec un monitoring automatique des logs.
\end{itemize}
\section{Réalisation des Modules Spécialisés}

\subsection{Implémentation de la Génération de PV}
Côté Frontend, nous avons implémenté une logique de calcul robuste pour transformer les notes brutes en un tableau de délibération officiel conforme au système tunisien.

\begin{lstlisting}[language=JavaScript, caption=Algorithme de calcul des moyennes (Frontend)]
function calculateStudentAverages(grades) {
    return grades.map(g => {
        const examWeight = 0.7;
        const ccWeight = 0.3;
        const total = (g.examGrade * examWeight) + (g.ccGrade * ccWeight);
        return {
            ...g,
            totalGrade: Math.round(total * 100) / 100,
            mention: getMention(total),
            isPassed: total >= 10
        };
    });
}
\end{lstlisting}

\subsection{Intégration de la Visioconférence (Daily.co)}
La gestion des sessions live repose sur la création dynamique de salles via l'API REST de Daily.co et leur affichage via un SDK React personnalisé.

\begin{lstlisting}[language=JavaScript, caption=Création d une session vidéo côté serveur]
async function createLiveSession(title) {
    const response = await axios.post('https://api.daily.co/v1/rooms', {
        name: `room-${Date.now()}`,
        properties: {
            enable_chat: true,
            start_video_off: false,
            exp: Math.round(Date.now() / 1000) + 7200 // Expire dans 2h
        }
    }, { headers: { Authorization: `Bearer ${PROCESS.env.DAILY_KEY}` } });
    return response.data;
}
\end{lstlisting}

\section{Réalisation des Interfaces Utilisateurs}
L'architecture frontend avec Next.js nous a permis de découper l'application en trois espaces de travail (\textit{Workspaces}) étanches, chacun ayant sa propre logique de routing et ses propres composants.

\subsection{L'Espace Administrateur : Dashboard de Pilotage}
L'administrateur dispose d'une vue macroscopique sur l'état de la plateforme. Nous avons utilisé des bibliothèques de graphiques (Recharts) pour afficher :
\begin{itemize}
    \item L'évolution du nombre d'inscriptions par mois.
    \item La consommation des tokens de l'API Gemini (monitoring des coûts).
    \item Le taux de réussite moyen aux quiz générés par l'IA.
\end{itemize}
\fbox{Placeholder: Capture d'écran - Panel Admin Dashboard}

\subsection{L'Espace Enseignant : Studio de Création Pédagogique}
L'interface enseignant est centrée sur la productivité. La gestion des chapitres se fait via une interface intuitive permettant de glisser-déposer des fichiers PDF.
\begin{itemize}
    \item \textbf{Flux de l'IA} : Lorsqu'un enseignant clique sur "Générer", un skeleton de chargement (\textit{Shimmer Effect}) s'affiche pendant que le serveur traite la demande Gemini. Le résultat est affiché dans un éditeur de texte riche (TipTap) pour modification.
    \item \textbf{Gestion des Live} : Un tableau liste les sessions passées et futures, avec accès direct aux statistiques de présence générées à partir des logs Daily.co.
\end{itemize}
\fbox{Placeholder: Capture d'écran - Interface de Génération IA Enseignant}

\subsection{L'Espace Étudiant : Environnement d'Étude Immersif}
L'espace étudiant privilégie la clarté et l'absence de distractions.
\begin{itemize}
    \item \textbf{Lecteur de Cours Intelligent} : Le PDF est affiché côte à côte avec le résumé généré par l'IA.
    \item \textbf{Module d'Examen} : Les quiz sont présentés sous forme d'une application de type "Quiz App", avec une barre de progression et une animation lors du passage d'une question à l'autre.
    \item \textbf{Chat Multimodal} : L'étudiant peut poser des questions au professeur ou à l'ensemble de la classe, avec support du format Markdown pour les extraits de code.
\end{itemize}
\fbox{Placeholder: Capture d'écran - Hub Étudiant et Salle de Vidéo}

\section{Implémentation du Service d'IA Locale (Ollama)}
Pour intégrer le modèle Gemma 3:4b, nous avons développé une classe \texttt{OllamaService}. Ce service communique avec l'instance locale via des requêtes HTTP sur le port 11434.

\begin{lstlisting}[language=JavaScript, caption=Implémentation du service Ollama pour l'IA locale]
import axios from 'axios';

class OllamaService {
  constructor() {
    this.baseURL = 'http://localhost:11434/api/generate';
  }

  async generateSummary(text) {
    const response = await axios.post(this.baseURL, {
      model: 'gemma3:4b',
      prompt: `Resume ce texte de maniere ludique : ${text}`,
      stream: false
    });
    return response.data.response;
  }
}
\end{lstlisting}

\section{Optimisation des Performances}
\subsection{Caching des Réponses IA}
La génération de contenu par un LLM peut être lente (de 5 à 15 secondes). Pour améliorer l'expérience utilisateur, nous avons mis en place un système de cache avec \textbf{Redis}.
\begin{itemize}
    \item Si un enseignant demande le résumé d'un cours déjà traité, le serveur renvoie la réponse instantanément depuis Redis.
    \item Les clés de cache sont basées sur le \texttt{SHA-256} du contenu du fichier PDF.
\end{itemize}

\subsection{Code Splitting et Lazy Loading (Frontend)}
Pour réduire le temps de chargement initial (\textit{First Contentful Paint}), nous utilisons le \textbf{Dynamic Import} de Next.js pour les composants lourds comme le lecteur PDF et la salle de visioconférence.
\begin{lstlisting}[language=JavaScript, caption=Exemple de chargement dynamique du composant Video]
const DailyVideo = dynamic(() => import('@/components/video/DailyVideo'), {
  ssr: false,
  loading: () => <Skeleton className="h-full w-full" />
});
\end{lstlisting}

\section{Gestion des Erreurs et Standardisation}
...

Pour offrir une expérience utilisateur fluide, nous avons standardisé la gestion des erreurs côté backend.

\subsection{Middleware de Gestion d'Erreurs Centralisé}
Toutes les exceptions sont interceptées par un middleware unique qui formate la réponse selon le standard suivant :
\begin{lstlisting}[language=JavaScript, caption=Structure d'une réponse d'erreur standard]
{
    "status": "error",
    "message": "Description lisible de l'erreur",
    "stack": "Trace (uniquement en developpement)"
}
\end{lstlisting}

\subsection{Journalisation (Logging)}
Nous utilisons la bibliothèque \textbf{Winston} pour journaliser les erreurs critiques dans des fichiers tournants sur le serveur, facilitant ainsi le débogage post-déploiement.

\section{Conclusion}
La phase de réalisation a permis de transformer les modèles conceptuels en une application robuste. L'intégration de l'IA Google Gemini s'est avérée être le moteur de différenciation d'eduGenius par rapport aux solutions classiques du marché.
